\documentclass[10pt,a4paper]{article}
\usepackage{amsmath,amssymb,amsthm,graphicx}
\usepackage{hyperref}
\usepackage{enumitem}
\usepackage{mathtools}
\title{An exponential-hierarchy oracle prevents almost-everywhere subexponential pseudodeterministic constructions}
\author{Shuichi Hirahara, Zhenjian Lu, Hanlin Ren}
\date{doi:10.4230/LIPIcs.CCC.2023.6; CC BY 4.0}
\begin{document}
\maketitle
\noindent\textit{Editorial note.} Complete author official proof faithfully transcribed from CC BY4 published PDF. Author mathematics and original labels retained. Standard cited Find1 and probability-class facts are retained without recursive reproof. No native author TeX was located through the official/ECCC/author publication routes.\par\medskip
\noindent\textit{Allowed standard prerequisites (editorial).}\begin{itemize}
\item Pseudodeterministic Find1 query lower bound and BPP containment in the second level. Exact original Theorems51 [Goldwasser-Impagliazzo-Pitassi-Santhanam2021] and23 [Sipser1983] retained as author-cited standard inputs
\end{itemize}\medskip
\paragraph{Definitions and notations (original Sections 1.1, 2.1).}
For a complexity class $\mathfrak C$, a complexity-theoretic statement is $\mathfrak C$-relativizing if it is true relative to every oracle $\mathcal O\in\mathfrak C$.
\paragraph{Definition 20.} A language $L$ is in $\mathsf{EXPH}$ if there is a constant $k$, an exponential bound $e(n)=2^{\operatorname{poly}(n)}$, and a $\operatorname{poly}(e(n))$-time algorithm $M$ such that, for every input $x\in\{0,1\}^*$,
\[x\in L\Longleftrightarrow\exists y_1\forall y_2\exists y_3\forall y_4\ldots\exists y_k,\ M(x,y_1,\ldots,y_k)=1,\]
where $y_i\in\{0,1\}^{e(n)}$ for each $i\in[k]$. Equivalently, a language is in $\mathsf{EXPH}$ if it can be decided in exponential ($2^{\operatorname{poly}(n)}$) time with access to a $\mathsf{PH}$ oracle.
\paragraph{Pseudodeterministic algorithms.} A pseudodeterministic algorithm is a probabilistic algorithm that returns a fixed output with high probability on any given input. We will also consider pseudodeterministic oracle (or advice-taking) algorithms. Such an algorithm will maintain its pseudodeterministic behavior (i.e., outputting a fixed answer with high probability for any input) when a correct oracle (or advice) is given.
\paragraph{Dense property (original Section 1.3.2).} A property $Q\subseteq\{0,1\}^*$ is dense if for every input length $n\in\mathbb N$, we have that $|Q\cap\{0,1\}^n|\geq2^n/\operatorname{poly}(n)$.
\paragraph{Theorem 23 [68].} For every oracle $\mathcal O$, $\mathsf{BPP}^{\mathcal O}\subseteq(\Sigma^p_2)^{\mathcal O}$.

\subsection*{5.2 Bounded-Relativization Barriers for Explicit Constructions}
\subsubsection*{5.2.1 Barriers for Almost-Everywhere Pseudodeterministic Constructions}
We show that $\mathsf{EXPH}$-relativizing techniques cannot prove almost-everywhere pseudodeterministic constructions, even if the construction algorithms are allowed $2^{o(n)}$ time. The underlying oracle uses the query complexity lower bounds proved by Goldwasser, Impagliazzo, Pitassi, and Santhanam [31] in a black-box fashion. (We thank Rahul Santhanam for pointing out that the query complexity lower bounds imply an $\mathsf{EXPH}$-computable oracle without almost-everywhere pseudodeterministic constructions in a black-box fashion.)
\paragraph{Definition 50.} A pseudodeterministic decision tree for a search problem $S$ is a distribution $\mathcal T$ over decision trees with the following property: For every input $x$, there is a canonical value $o$ such that with probability at least $2/3$, $\mathcal T(x)=o$. Let $\operatorname{psPdt}(S)$ denote the minimum depth of any pseudodeterministic decision tree for $S$.
Consider the following search problem denoted as $\mathsf{Find1}$. The input is a string $x\in\{0,1\}^N$ where it is guaranteed that there are at least $N/2$ bits in $x$ that are equal to $1$. The problem is to find some $i\in[N]$ such that $x_i=1$. The following lower bound on the pseudodeterministic query complexity of $\mathsf{Find1}$ was proved in [31]:
\paragraph{Theorem 51 [31].} $\operatorname{psPdt}(\mathsf{Find1})=\Omega(\sqrt N)$.
We use this result as a black box and construct an oracle $\mathcal O\in\mathsf{EXPH}$ without almost-everywhere pseudodeterministic constructions running in $2^{o(n)}$ time:

\paragraph{Theorem 52.} There is an oracle $\mathcal O\in\mathsf{EXPH}$ and a dense property $Q\in\mathsf P^{\mathcal O}$ such that the following holds. For every randomized (and purportedly pseudodeterministic) algorithm $A$ that runs in $2^{o(n)}$ time with oracle access to $\mathcal O$ and every infinite sequence of outputs $\{x_n\}_{n\in\mathbb N}$ where each $x_n\in Q\cap\{0,1\}^n$, there are infinitely many input lengths $n\in\mathbb N$ such that
\[\Pr[A(1^n)=x_n]\leq3/4.\]

\begin{proof}
Let $M_1,M_2,\ldots$ be an enumeration of probabilistic machines running in $2^{0.1n}$ time. Let $n_1$ be a large enough constant and define the sequence $\{n_i\}$ where $n_i=2^{n_{i-1}}$ for each $i\geq2$. We let $Q=\mathcal O$ itself be the dense property without pseudodeterministic construction algorithms. In particular, we want that for every $i\in\mathbb N$, $M_i(1^{n_i})$ fails to generate a canonical string $x\in\mathcal O\cap\{0,1\}^{n_i}$. Naturally, our oracle construction proceeds in stages, where for each $i\in\mathbb N$, the $n_i$-th slice of $\mathcal O$ is constructed by carefully diagonalizing against $M_i$; on the other hand, if $n\in\mathbb N$ is not in the sequence $\{n_i\}$, then we simply let $\mathcal O$ accept every string of length $n$.

Now we show how to construct the $n_i$-th slice of $\mathcal O$ in the $i$-th stage. The idea is simple: construct the (purportedly pseudodeterministic) decision tree $\mathcal T$ corresponding to $M_i$, find an exponential-length input $tt$ of $\mathsf{Find1}$ on which $\mathcal T$ fails, and let the truth table of the $n_i$-th slice of $\mathcal O$ be $tt$. More precisely:
\begin{itemize}
\item \textbf{Converting $M_i$ into a distribution of decision trees.} We define the distribution of decision trees $\mathcal T$. The input to $\mathsf{Find1}$ has length $2^{n_i}$ and is considered as the truth table of the $n_i$-th slice of $\mathcal O$. Note that since $2^{0.1n_i}\ll n_{i+1}$, for every $n'\leq2^{0.1n_i}$ such that $n'\neq n_i$, the $n'$-th slice of $\mathcal O$ are fixed. To sample a (deterministic) decision tree $T$ from $\mathcal T$, sample a sequence of random coins fed to $M_i$ and simulate $M_i$ on the oracle $\mathcal O$. Whenever $M_i$ makes a query $\mathcal O(x)$, if $|x|=n_i$, then the decision tree asks the $x$-th bit of our input; otherwise we return the already-fixed value of $\mathcal O(x)$.
\item \textbf{Invoking the lower bound.} Note that $\mathcal T$ only makes $2^{0.1n_i}<o(\sqrt{2^{n_i}})$ queries. By Theorem 51, there exists an input $tt\in\{0,1\}^{2^{n_i}}$ such that $\mathcal T$ fails to solve $\mathsf{Find1}$ on input $tt$ pseudodeterministically. That is, for every valid output $x$ such that $tt_x=1$, $\Pr[\mathcal T(tt)=x]<2/3$. Let the truth table of the $n_i$-th slice of $\mathcal O$ to be such a $tt$.
\end{itemize}
It is easy to see that for every probabilistic machine $A$ that runs in $2^{o(n)}$ time with oracle access to $\mathcal O$, and every infinite sequence of outputs $\{x_n\}_{n\in\mathbb N}$ where each $x_n\in\mathcal O\cap\{0,1\}^n$, for every $i$ such that $A=M_i$, the probability that $A(1^{n_i})$ outputs $x_{n_i}$ is at most $3/4$.

\textbf{Complexity of $\mathcal O$.} Let $\mathcal O_{\leq n_{i-1}}$ denote the description of the oracle $\mathcal O$ up to the $n_{i-1}$-th slice (note that $\mathcal O_{\leq n_{i-1}}$ can be described in $\operatorname{poly}(2^{n_{i-1}})\leq\operatorname{poly}(n_i)$ bits), and $x\in\{0,1\}^{n_i}$. Note that $\mathcal O_{\leq n_{i-1}}$ and $tt$ together defines the oracle $\mathcal O$ up to input length $n_{i+1}-1$, so the behavior of $M_i^{\mathcal O}$ is completely determined by $M_i$, $\mathcal O_{\leq n_{i-1}}$, and $tt$. Let $\mathsf{ProbEst}(\mathcal O_{\leq n_{i-1}},tt,M_i,x)$ be an oracle such that
\begin{align*}
\Pr[M_i^{\mathcal O}(1^{n_i})=x]>3/4&\Longrightarrow\mathsf{ProbEst}(\mathcal O_{\leq n_{i-1}},tt,M_i,x)=1,\\
\Pr[M_i^{\mathcal O}(1^{n_i})=x]<2/3&\Longrightarrow\mathsf{ProbEst}(\mathcal O_{\leq n_{i-1}},tt,M_i,x)=0.
\end{align*}
By Theorem 23, the oracle $\mathsf{ProbEst}$ can be implemented in $\mathsf{DTIME}[2^{O(n_i)}]^{\mathsf{PH}}$. Consider the following algorithm with oracle access to $\mathsf{ProbEst}$ that prints a truth table $tt$. The algorithm maintains the prefix of a truth table which is initially the empty string, and extends this prefix bit by bit. Suppose that we have a prefix of length $\ell$, denoted as $tt'\in\{0,1\}^{\ell}$. To fix the $(\ell+1)$-st bit of $tt'$, we check if there exists a truth table $tt$ such that (1) for every $x\in\{0,1\}^{n_i}$, $\mathsf{ProbEst}(\mathcal O_{\leq n_{i-1}},tt,M_i,x)=0$, and (2) $tt'\circ0$ ($tt'$ concatenated with a bit $0$) is a prefix of $tt$. If there is such a $tt$, then we append $0$ to the end of $tt'$; otherwise we append $1$ to the end of $tt'$.

It is clear that the algorithm always outputs a truth table $tt$ on which $M_i$ fails. Since $\mathsf{ProbEst}$ runs in $\mathsf{DTIME}[2^{O(n_i)}]^{\mathsf{PH}}$, the whole algorithm also runs in $\mathsf{DTIME}[2^{O(n_i)}]^{\mathsf{PH}}$. Using this algorithm, it is easy to construct the oracle $\mathcal O$ in $\mathsf{EXPH}$.
\end{proof}

\begin{thebibliography}{99}
\bibitem[31]{GIPS} Shafi Goldwasser, Russell Impagliazzo, Toniann Pitassi, and Rahul Santhanam. On the pseudo-deterministic query complexity of NP search problems. CCC2021, LIPIcs200,36:1--36:22. doi:10.4230/LIPIcs.CCC.2021.36.
\bibitem[33]{H18} Shuichi Hirahara. Non-black-box worst-case to average-case reductions within NP. FOCS2018,247--258. doi:10.1109/FOCS.2018.00032.
\bibitem[34]{H20} Shuichi Hirahara. Non-disjoint promise problems from meta-computational view of pseudorandom generator constructions. CCC2020,LIPIcs169,20:1--20:47. doi:10.4230/LIPIcs.CCC.2020.20.
\bibitem[68]{Sipser} Michael Sipser. A complexity theoretic approach to randomness. STOC1983,330--335. doi:10.1145/800061.808762.
\end{thebibliography}
\end{document}
